\begin{circuitikz}[
    american,
    voltage dir=RP,
    >=latex,
    %>={Latex[length=5mm, width=2mm]},
    alt={Diagram showing a ladder logic program containing a Retentive Timer On (RTO) instruction and a Reset (RES) instruction across three rungs. The top rung contains a normally open contact Auto_Mode connected to an RTO timer block named AutoTimer with a Preset (PRE) value of 28800 and an Accumulated (ACC) value of 0. The bottom rung contains a normally open contact ShiftChange connected to a Reset instruction block labeled RES targeting AutoTimer. Green power flow highlighting extends from the left vertical rail through the ShiftChange contact and across the entire bottom rung through the RES instruction, indicating that the ShiftChange condition is true, executing the reset and holding the AutoTimer accumulated value at 0.}
    ]

    % Define the number of rungs in the diagram
    \def\numRungs{3}

    % Calculate the length of the vertical power rails dynamically based on number of rungs
    \pgfmathsetmacro{\railLength}{-2 * \numRungs}

    %======================================================================================================
    % Component Grid
    % Spacing = 3x, 2y
    %======================================================================================================
    \foreach \i in {0,...,5} {
        \foreach \j in {0,...,20} { % Change to 0,...,19 if you need 20 individual Y points down to -40
        
        % Calculate spatial positions
        \pgfmathsetmacro{\px}{\i * 3}
        \pgfmathsetmacro{\py}{-1 - (\j * 2)}
        
        % Name coordinate using loop indices: (C-column-row) e.g., (C-0-0), (C-5-10)
        \coordinate (C-\i-\j) at (\px, \py);
        
        % OPTIONAL: Uncomment the 2 lines below to display grid labels while building
        % \fill[red] (C-\i-\j) circle (1.5pt);
        %\node[anchor=south west, scale=0.5, gray] at (C-\i-\j) {(\i,\j)};
        }
    }

    %======================================================================================================
    % Foreground Layer
    % Only contains the vertical power rails
    % Length of the rails is automatic, calculated based on the number of rungs defined above
    %======================================================================================================
    \begin{pgfonlayer}{ft}

        %======================================================================================================
        % Scope sets options for both lines - BakerBlack, very thick
        %======================================================================================================
        \begin{scope}[
            draw=BakerBlack,
            text=BakerBlack,
            color=BakerBlack,
            font=\small,
            ultra thick
        ]
        
            \draw (0,0) -- (0,\railLength);

            \draw (15,0) -- (15,\railLength);

        \end{scope}
    \end{pgfonlayer}

    %======================================================================================================
    % Main Layer
    % All actual logic goes here
    %======================================================================================================
    \begin{pgfonlayer}{main}
        
        %======================================================================================================
        % Scope sets options for both lines - BakerBlack for draw and text
        %======================================================================================================
        \begin{scope}[
            draw=BakerBlack,
            text=BakerBlack,
            color=BakerBlack,
            font=\small
        ]

            %=================================================================
            % Rung 1
            %=================================================================
            \draw(C-0-0) to[C, l={Auto\_Mode}] (C-1-0) -- (C-3-0) pic (T1) {TON};

            \draw (T1-OUT) -- (C-5-0);

            \node[anchor=south] at (T1-TOP) {RTO};
            \node[anchor=north] at (T1-TOP) {AutoTimer};
            \node[anchor=west] at (T1-PRE) {28800};
            \node[anchor=west] at (T1-ACC) {0};

            %=================================================================
            % Rung 3 (timer takes 2 rungs)
            %=================================================================
            \draw (C-0-2) to[C, l={ShiftChange}] (C-1-2) -- (C-4-2)
                node[
                    draw=BakerBlack,
                    minimum width=1cm, 
                    minimum height=0.5cm, 
                    anchor=west, 
                    rounded corners=0.25cm
                ](res)
                {RES};

            \node[above] at (res.north){AutoTimer};

            \draw (res.east) -- (C-5-2);

        \end{scope}
    \end{pgfonlayer}

    %=================================================================
    % Background layer
    %=================================================================
    \begin{pgfonlayer}{bg}
        
        %==================================================================
        % Highlight Scope
        % This layer exists for highlighting true logic, the scope sets the
        % highlight options without having to repeat them every \draw
        % draw=UpNorthGreen, line width=10pt, opacity=0.4
        %==================================================================
        \begin{scope}[
            draw=UpNorthGreen,
            line width=10pt,
            opacity=0.4
        ]

            %==================================================================
            % Left vertical rung - leave this alone, it should ALWAYS be green
            % It automatically draws itself to the length defined at the top
            %==================================================================
            \draw (0,0) -- (0,\railLength);
            
            %=================================================================
            % Rung 3
            %=================================================================
            \draw (C-0-2) -- (C-5-2);

        \end{scope}
    \end{pgfonlayer}
    
\end{circuitikz}